% source: luadraw-v3.5/luadraw/doc/src/body-en/luadraw-doc3d-en-section-7-Matrix-Calculus-Transformations-Some-Mathematical-Functions.tex (demo: A curve on a cylinder)
\begin{luadraw}{name=curve_on_cylinder}
local ld = luadraw
local pt3d = ld.pt3d
local Origin, vecI, vecJ, vecK, M, Mc = pt3d.Origin, pt3d.vecI, pt3d.vecJ, pt3d.vecK, pt3d.M, pt3d.Mc

local g = ld.graph3d:new{adjust2d=true,bbox=false,size={10,10}, viewdir="central"}
g:Labelsize("footnotesize")
ld.Hiddenlines = true; ld.Hiddenlinestyle = "dashed"
local curve_on_cylinder = function(curve,cylinder)
-- curve is a 3D polyline on a cylinder,
-- cylinder = {A,r,V,B}
    local A,r,V,B = table.unpack(cylinder)
    if B == nil then B = V; V = B-A end
    local U = B-A
    local visible_function
    if ld.projection_mode == "central" then
        visible_function = function(N)
            local I = ld.dproj3d(N,{A,U})
            local M1, M2 = ld.interCS({I,r,U},{ (I+ld.camera)/2, pt3d.abs(I-ld.camera)/2})
            local alpha = pt3d.angle3d(M1-I,ld.camera-I)
            return pt3d.angle3d(N-I,ld.camera-I) <= alpha
        end
    else
        visible_function = function(N)
            local I = ld.dproj3d(N,{A,U})
            return (pt3d.dot(N-I,g.Normal) >= 0)
        end
    end
    return ld.split_points_by_visibility(curve,visible_function)
end
-- test
local A, r, B = -5*vecJ, 4, 5*vecJ -- cylinder
local p = function(t) return Mc(r,t,t/3) end
local Curve = ld.rotate3d( ld.parametric3d(p,-4*math.pi,4*math.pi),90,{Origin,vecI})
local Vi, Hi = curve_on_cylinder(Curve,{A,r,B})
local curve_color = "DarkGreen"
g:Dboxaxes3d({grid=true,gridcolor="gray",fillcolor="LightGray"})
g:Dcylinder(A,r,B,{color="orange"})
g:Dpolyline3d(Vi,curve_color)
g:Dpolyline3d(Hi,curve_color..","..ld.Hiddenlinestyle)
g:Show()
\end{luadraw}
